\section{Introduction}
\label{sec:intro}

The ability to merge multiple specialized deep neural networks into a single multi-task model without additional expensive training has emerged as a promising frontier in deep learning. Traditional model merging techniques, such as Task Arithmetic \cite{ilharco2022editing} and TIES-Merging \cite{yadav2023ties}, operate under a fundamental assumption: that the merged parameters must form a \emph{static, input-independent} point in the Euclidean weight-space. Even test-time adaptation methods like AdaMerging \cite{liang2023adamerging}, which adjust layer-wise merging coefficients dynamically, ultimately seek a static consensus weight matrix per task or a uniform compromise across tasks.

While static merging works reasonably well for over-parameterized models (e.g., large-scale language models) or low-conflict domains, it breaks down catastrophically when applied to compact networks or highly diverse, conflicting tasks. On a compact backbone like a 5.7M-parameter Vision Transformer ($\mathtt{vit\_tiny\_patch16\_224}$), standard static merging results in what we define as \emph{catastrophic representational collapse}. When attempting to merge expert models trained on vastly different visual distributions (e.g., MNIST, FashionMNIST, CIFAR-10, and SVHN) which have been trained to high convergence (with an average specialized performance ceiling of $70.52\%$), the conflicting parameter shifts destroy the delicate representational pathways. Under uniform weight averaging, the joint multi-task performance collapses to $49.35\%$, with FashionMNIST plummeting to $44.50\%$ and SVHN to $21.90\%$.

In this work, we challenge this static-parameter assumption. Drawing inspiration from quantum mechanics, we propose a radical and non-incremental paradigm shift: \textbf{Quantum Wavefunction Superposition Merging (QWS-Merge)}. Rather than forcing multiple expert models to compromise on a single, static point in weight-space, we model each task expert as a task eigenstate $|\psi_k\rangle$ in a parameter Hilbert space. Merging is thus formulated as a coherent, quantum-like superposition of these states, which dynamically ``collapses'' (is measured) to a localized classical weight configuration based on the phase-overlap of the incoming input features.

Specifically, we project the global feature representation of each input batch onto a low-dimensional unit sphere, representing the input phase state. We then optimize a set of layer-wise trainable phase basis vectors, scaling amplitudes, and phase biases on a tiny offline validation set (only 16 samples per task, total 64 calibration images). At inference time, the dynamic probability amplitude (merging coefficient) for each expert is determined on-the-fly by wave-like phase interference between the input phase state and the learned phase-basis vectors of the layers. These sample-level amplitudes are averaged across the batch to collapse the wavefunction into a classical, batch-level weight configuration. By routing representations dynamically through task-specific activation pathways, QWS-Merge completely resolves the spatial weight-cancellation problem that plagues linear averaging, allowing a single compact model to support highly conflicting visual domains simultaneously.

Crucially, our revised study introduces a classical dynamic routing baseline, the \emph{Linear Router}, operating under an identical optimization budget and comparable parameter footprint. We make a striking empirical discovery: while classical linear routing achieves a slightly higher overall joint mean accuracy of $61.23\%$, it is highly prone to unconstrained parameter-space collapse and overfitting when confronted with extreme task conflicts. On the highly challenging SVHN dataset, the Linear Router collapses to a meager $15.30\%$ accuracy. In stark contrast, QWS-Merge's cosine-based wave phase-interference projections act as a highly regularized, bounded subspace, maintaining a high performance of \textbf{31.60\%} on SVHN (preserving $91.5\%$ of its specialized expert capacity). Furthermore, we conduct the first systematic analysis of batch size and task heterogeneity on dynamic merging, exposing how mixed-task batches lead to ``heterogeneity collapse'' at larger batch sizes, providing a transparent and scientifically honest benchmark for future research.

Our contributions are summarized as follows:
\begin{itemize}
    \item \textbf{A Visionary Paradigm Shift:} We challenge the static-parameter assumption of model merging and propose a quantum-inspired formulation where experts are task eigenstates in a Hilbert space, superposed and collapsed dynamically via input phase-coherence.
    \item \textbf{High-Conflict Model Merging Solution:} We demonstrate that QWS-Merge completely resolves catastrophic representational collapse on compact backbones in high-conflict scenarios, a regime where standard static baselines suffer from severe parameter interference.
    \item \textbf{Wave-Like Subspace Regularization:} We prove that QWS-Merge's non-monotonic cosine phase projections provide robust regularization. Under extreme task conflict (SVHN), our method outscores a classical soft routing baseline by a margin of $+16.30\%$, preventing parameter-space collapse.
    \item \textbf{Transparent Heterogeneity Benchmark:} We present a systematic investigation into test batch size and task heterogeneity, illustrating the trade-offs of dynamic parameter-space routers on mixed-task streams.
    \item \textbf{Extreme Resource Efficiency:} We show that our method optimizes only 336 parameters on a tiny 64-sample offline validation set, completely bypassing the Overfitting-Optimizer Paradox associated with high-dimensional parameter adaptation.
\end{itemize}
